\chapter{Exploration versus Exploitation}

\chapterepigraph{%
  The explorer who never settles learns everything about everything\\
  and masters nothing.\\
  The expert who never explores masters one world\\
  and misses all the others.\\
  Wisdom is the calibration between them.
}{Flyxion, \textit{Semantic Infrastructure}}

\sidenote{App.\,N §N.16\\on the exploration-\\exploitation tradeoff\\in curricula}

The exploration-exploitation tradeoff is one of the deepest recurring
problems in learning, decision theory, and evolutionary biology. In the
knowledge infrastructure framework, it takes a precise form: given a
finite maintenance budget, how should an agent allocate effort between
exploring new regions of the knowledge landscape (building new connections)
and deepening existing regions (strengthening established structures)?

\section{Discovery}

Discovery is the acquisition of genuinely new constraints: knowledge
that was not previously representable in the agent's infrastructure.
It requires leaving the current high-accessibility regions and venturing
into lower-accessibility territory where new structure may be found.

\begin{definition}{Epistemic Exploration}{epistemic-exploration}
  An agent is in \emph{epistemic exploration mode} when it preferentially
  visits concepts $k$ with high uncertainty $\sigma_K(k)$ (unknown
  territory) over concepts with high current accessibility $S_K(k)$
  (familiar territory). The exploration bonus is:
  \[
    \text{UCB}(k) = S_K(k) + \beta \cdot \sigma_K(k)
  \]
  (Upper Confidence Bound), where $\beta$ controls the
  exploration-exploitation tradeoff.
\end{definition}

\section{Practice}

Practice is the reinforcement of existing structures: visiting known
concepts to reduce their decay rate, strengthen their connections to
neighboring concepts, and increase the precision of the associated
constraint modifications.

\begin{definition}{Deliberate Practice}{deliberate-practice}
  \emph{Deliberate practice} is practice at the boundary of current
  competence: visiting concepts just beyond the current high-accessibility
  region, where the constraint modifications are not yet robust.
  Formally, it targets concepts with moderate current accessibility
  $S_K(k)$ (accessible but not automatic) and high rate-of-improvement
  $\partial S_K(k)/\partial t$ under practice.
\end{definition}

This formalizes the pedagogical observation that learning happens at
the edge of competence. Too easy (within the high-accessibility region)
produces no new constraint modifications; too hard (far outside it)
produces confusion and anxiety. The productive zone is the boundary.

\section{Mastery and Novelty}

Mastery is the state where a concept's accessibility is so high that
it requires minimal maintenance and can be used automatically in other
contexts. A master has high-$S_K$ coverage of their domain and can
freely combine concepts from different parts of it.

Novelty is the state where a concept has been recently acquired and is
not yet well-integrated into the knowledge infrastructure. Novel concepts
have high potential (they may open new regions) but high uncertainty
(their full implications are not yet clear).

\begin{namedthm}{The Exploration-Exploitation Theorem for Learning}
  An agent with finite learning time $T$ and knowledge infrastructure
  $(\mathcal{C}, S_K)$ should allocate time $T_\text{explore}$ to
  exploration and $T_\text{exploit}$ to exploitation in proportion to:
  \[
    \frac{T_\text{explore}}{T_\text{exploit}} =
    \frac{\max_k \sigma_K(k)}{S_K(k^*)_\text{avg}}
  \]
  — the ratio of maximum uncertainty to average current accessibility.
  Early in learning (high uncertainty, low accessibility), explore more.
  Late in learning (low uncertainty, high accessibility), exploit more.
\end{namedthm}

\begin{exercise}
  Track your own learning behavior for one week. For each study or
  reading session, classify it as primarily exploration (new territory)
  or exploitation (reinforcing existing knowledge). Are you in the
  right regime for your current stage of learning? What would the
  Exploration-Exploitation Theorem prescribe?
\end{exercise}

\begin{exercise}
  Scientific fields go through cycles of exploration (paradigm shifts,
  new methods) and exploitation (normal science, application of
  established methods). Identify a field you know well and describe
  its current exploration-exploitation balance. Is it in the right
  regime? What would shift it?
\end{exercise}
